\begin{table}[htbp]
\centering\small
\caption{Coverage and conditional runtime by family.}
\label{tab:family-runtime}
\input{\SubmissionSummaryDir/family_runtime.tex}
\tablenote{Each row includes the three graph sizes. SAT OPT and ILP OPT
count instances optimal within budget in all three repetitions; Joint
is their intersection. Both timing columns use only that same Joint
subset: first the three-run median for each instance, then the median
across instances. The ALL timings cover 74 of 117 instances and exclude
the other 43. A dash means the subset is empty. ER densities are
separated; ALL is computed across instances, not group medians.}
\end{table}

Runtime comparisons depend on the graph family and parameter pair.
On the jointly all-three-OPT subsets, CaDiCaL has lower medians for
trees, grids and BA graphs. The direction reverses for ER $p=0.30$:
the SAT and ILP medians are 7.196 and 0.470 seconds for $(1,1)$, and
4.055 and 0.811 seconds for $(2,1)$. Each of these timing comparisons
contains only two instances out of nine: the 20-vertex graphs with
seeds 0 and 2. No 40- or 60-vertex instance enters either timing subset.
ILP also has greater
all-three-OPT coverage in these groups: 8 versus 2 and 8 versus 3,
respectively. For ER $p=0.30$, $(3,2)$, the joint subset is empty,
so no conditional timing median is reported.
Trees, grids and BA graphs account for 61 of the 74 joint instances;
the two ER densities contribute nine and four, respectively.
The pooled medians therefore depend on both cohort composition and
selection into the jointly solved subset. They do not establish an
overall SAT speed advantage or a general ordering of graph difficulty.
